\documentclass{article}
\usepackage{listings}
\usepackage{xcolor}

\lstset{
    basicstyle=\ttfamily\small,
    keywordstyle=\color{blue},
    commentstyle=\color{gray},
    stringstyle=\color{red},
    numbers=left,
    numberstyle=\tiny,
    frame=single,
    breaklines=true
}

\title{Code Listings Example}
\author{latex-rs}

\begin{document}
\maketitle

\section{Rust Code}

\begin{lstlisting}[language=Rust]
fn main() {
    println!("Hello, World!");
    
    let numbers = vec![1, 2, 3, 4, 5];
    let sum: i32 = numbers.iter().sum();
    
    println!("Sum: {}", sum);
}
\end{lstlisting}

\section{Python Code}

\begin{lstlisting}[language=Python]
def fibonacci(n):
    if n <= 1:
        return n
    return fibonacci(n-1) + fibonacci(n-2)

for i in range(10):
    print(f"F({i}) = {fibonacci(i)}")
\end{lstlisting}

\end{document}
